\section{CPU 架构与模块设计}\label{sec:architecture}

\subsection{总体架构}

CPU 采用 IFU、IDU、EXU、LSU、WBU 五级顺序单发射流水线。级间寄存器采用 Decoupled
ready/valid 语义：后级不能接收时，前级保持当前 payload；操作真正完成时才发生 \code{fire}。
控制流错误通过 epoch 翻转使年轻指令局部失效，并由注册化重定向邮箱把正确 PC 送回 IFU。

\begin{figure}[H]\centering\resizebox{\linewidth}{!}{%
\begin{tikzpicture}[s/.style={draw,rounded corners,minimum width=2.15cm,minimum height=1cm,align=center,fill=blue!6},q/.style={draw,minimum width=1.05cm,minimum height=0.65cm,fill=gray!8},a/.style={-Stealth,thick},font=\small]
\node[s] (i0) {IFU}; \node[q,right=0.25cm of i0] (q0) {IF/ID}; \node[s,right=0.25cm of q0] (i1) {IDU};
\node[q,right=0.25cm of i1] (q1) {ID/EX}; \node[s,right=0.25cm of q1] (i2) {EXU}; \node[q,right=0.25cm of i2] (q2) {EX/LS};
\node[s,right=0.25cm of q2] (i3) {LSU}; \node[q,right=0.25cm of i3] (q3) {LS/WB}; \node[s,right=0.25cm of q3] (i4) {WBU};
\foreach \x/\y in {i0/q0,q0/i1,i1/q1,q1/i2,i2/q2,q2/i3,i3/q3,q3/i4}{\draw[a] (\x)--(\y);}
\draw[a,red!70!black] (i2.south) -- ++(0,-0.65) -| node[pos=0.2,below]{redirect + epoch} (i0.south);
\draw[a,green!50!black] (i4.north) -- ++(0,0.65) -| node[pos=0.25,above]{写回/旁路状态} (i1.north);
\end{tikzpicture}}\caption{五级流水、反压与重定向关系}\end{figure}


\subsection{IFU 取指模块}

IFU 将 PC、预测信息和流水 epoch 与读取到的指令组合为 \code{FetchedInst}。它能够在消费当前取指响应的
同一周期接受下一 PC 并发出下一请求，从而隐藏片上 IROM 的同步读延迟。若存储或 IDU 反压，则内部
pending 状态保持请求及其预测元信息，避免地址与预测结果错配。

\begin{signalstable}{IFU 接口信号说明}
pc&输入&32 位 Decoupled&\code{valid/ready}&接收待取指 PC\\
predNext&输入&PredBundle&与 PC 同周期&该 PC 对应的预测命中、方向和下一地址\\
epoch&输入&Bool&随 PC 采样&标记当前有效控制流世代\\
mem&输出&SimpleBus master&请求/响应握手&访问 IROM\\
out&输出&FetchedInst Decoupled&\code{valid/ready}&向 IDU 输出指令、PC、预测和 epoch\\
\end{signalstable}

\subsection{IDU 译码与相关处理模块}

IDU 完成指令格式/类型译码、立即数形成、寄存器读取和相关判定。它分别观察 EXU、LSU、WBU 的生产者
状态：目的寄存器匹配且数据已就绪时旁路；匹配但数据无效时停顿。冻结版本进一步把相邻快速结果按
8 个四位组编码，只允许普通整数、分支和 store-data 等短路径消费者使用；地址、M/B、自定义指令和
CSR 等消费者走注册化或保守停顿路径，以限制选择信号扇出。

\begin{signalstable}{IDU 接口信号说明}
in&输入&FetchedInst Decoupled&\code{valid/ready}&接收 IFU 结果\\
rvec&输出/输入&2 读端口 GPR&组合读&给出 rs1/rs2 地址并取得寄存器值\\
csrRead&输出/输入&CSR 单读口&组合读&读取当前 CSR 数据\\
wrBackInfo&输入&三级前递信息组&生产者有效性&判定旁路或停顿\\
lateLoadProducer&输入&延迟加载生产者信息&专用相关协议&标识可延后到 EXU 解析的 load 生产者\\
pipelineFlush&输入&Bool&高有效&丢弃错误 epoch 指令\\
out&输出&DecodedInst Decoupled&\code{valid/ready}&输出操作数、控制元信息和结果类别\\
\end{signalstable}

\begin{figure}[H]\centering
\begin{tikzpicture}[b/.style={draw,rounded corners,align=center,minimum width=2.5cm,minimum height=0.75cm,fill=blue!5},d/.style={draw,diamond,aspect=2,align=center,fill=orange!6},a/.style={-Stealth,thick},font=\small]
\node[b] (src) {源寄存器号}; \node[b,below=0.7cm of src] (p) {EXU/LSU/WBU 生产者优先级}; \node[d,below=0.8cm of p] (m) {目的寄存器匹配？};
\node[b,below left=0.9cm and 1.4cm of m] (g) {使用 GPR 数据}; \node[d,below right=0.9cm and 1.4cm of m] (v) {数据有效？};
\node[b,below left=0.9cm and 1.2cm of v] (s) {停顿}; \node[b,below right=0.9cm and 1.2cm of v] (f) {按消费者选择旁路};
\draw[a] (src)--(p);\draw[a](p)--(m);\draw[a](m)--node[above]{否}(g);\draw[a](m)--node[above]{是}(v);\draw[a](v)--node[above]{否}(s);\draw[a](v)--node[above]{是}(f);
\end{tikzpicture}\caption{IDU 相关判定与旁路流程}\end{figure}



\subsection{EXU 分区执行模块}

EXU 同时承担普通算术、分支裁决、地址请求、多周期整数扩展、CSR 操作和自定义指令执行。为避免所有
结果汇入同一宽 Mux，冻结版本把写回语义划分为 \code{fastInt}、\code{direct}、
\code{longArithmetic}、\code{accelerator} 和 \code{load} 五类。普通整数进入独立快速 ALU；
乘除法、迭代 B 扩展和短 B 运算进入特殊执行簇；领域单元使用独立结果通路；load 最终由存储响应产生。

\begin{signalstable}{EXU 接口信号说明}
in&输入&DecodedInst Decoupled&\code{valid/ready}&接收译码结果及消费者类型\\
previousStageFwd&输入&FastForwardInfo&局部组合选择&只供允许的相邻快速消费者使用\\
lateLoadLSU/WBU&输入&LateLoadInfo&有效/数据有效&在 EXU 注册边界解析受限 load-use\\
dcache&双向&DCache 查询/更新组&查询与更新协议&命中、同步/异步数据及 store 变更\\
memReq/memResp&双向&Decoupled/Valid&请求/响应&普通访存和自定义单元访存\\
fwd&输出&ForwardInfo&生产者状态&向 IDU 广告目的寄存器及数据有效性\\
out&输出&LSUInput Decoupled&\code{valid/ready}&输出分区写回结果、访存和重定向元信息\\
branch* / btbUpdate*&输出&控制流更新组&随完成采样&提供真实方向、目标和表项类型\\
\end{signalstable}

\begin{figure}[H]\centering
\begin{tikzpicture}[u/.style={draw,rounded corners=1.2pt,align=center,minimum width=2.55cm,minimum height=0.78cm},a/.style={-Stealth,thick},font=\small]
\node[u,fill=blue!6,minimum height=5.25cm] (d) {DecodedInst\\按 ResultKind 分发};
\node[u,fill=green!7,right=1.35cm of d,yshift=2.25cm] (fi) {快速整数通路};
\node[u,fill=green!7,right=1.35cm of d,yshift=1.12cm] (di) {直接结果通路};
\node[u,fill=orange!8,right=1.35cm of d] (la) {长算术通路};
\node[u,fill=purple!7,right=1.35cm of d,yshift=-1.12cm] (ac) {硬件加速通路};
\node[u,fill=cyan!7,right=1.35cm of d,yshift=-2.25cm] (ld) {加载结果通路};
\node[u,fill=gray!10,right=2.35cm of la,minimum height=5.25cm] (wb) {按 ResultKind\\局部汇聚与写回};
\draw[a] ([yshift=2.25cm]d.east)--(fi.west);
\draw[a] ([yshift=1.12cm]d.east)--(di.west);
\draw[a] (d.east)--(la.west);
\draw[a] ([yshift=-1.12cm]d.east)--(ac.west);
\draw[a] ([yshift=-2.25cm]d.east)--(ld.west);
\draw[a] (fi.east)--([yshift=2.25cm]wb.west);
\draw[a] (di.east)--([yshift=1.12cm]wb.west);
\draw[a] (la.east)--(wb.west);
\draw[a] (ac.east)--([yshift=-1.12cm]wb.west);
\draw[a] (ld.east)--([yshift=-2.25cm]wb.west);
\end{tikzpicture}\caption{EXU 与写回结果通路分区}\end{figure}


\subsection{LSU 与 WBU}

LSU 保存访存类型、地址低位、缓存命中和分区写回信息，并在同步缓存数据到达后形成 LSU 级旁路。
WBU 对 load 返回字按 \code{func3} 和地址偏移完成字节/半字抽取与符号扩展，随后写入 GPR；其他结果
按结果类别选择后顺序提交。store 的缓存更新在 EXU 请求侧发生，较早 load 的回填则由 store epoch
阻止覆盖更新后的缓存状态。

\begin{signalstable}{LSU 接口信号说明}
in&输入&LSUInput Decoupled&\code{valid/ready}&接收执行结果和访存元信息\\
dcacheReadData&输入&UInt(32)&同步 BRAM 返回&取得上一周期查询的数据\\
fwd&输出&ForwardInfo&生产者状态&向 IDU 提供 LSU 级普通结果状态\\
out&输出&WBUInput Decoupled&\code{valid/ready}&向 WBU 传递统一提交信息\\
\end{signalstable}

\begin{signalstable}{WBU 接口信号说明}
in&输入&WBUInput Decoupled&\code{valid/ready}&接收 LSU 输出\\
memResp&输入&Valid[UInt(32)]&与 load 对齐&接收 DRAM 返回数据\\
dcacheUpdate*&输出&地址/数据/掩码&提交时回填&将有效 load 结果写入缓存\\
gpr&输出&GPR WriteIO&最终提交&写入目的通用寄存器\\
\end{signalstable}

\subsection{GPR 与 CSR}

通用寄存器堆使用两份 32$\times$32 分布式存储器形成双异步读、单同步写结构；写回同址时使用显式
write-through，\code{x0} 始终返回零。CSR 单元实现 \code{mstatus}、\code{mtvec}、\code{mepc}、
\code{mcause}、\code{mcycle} 等机器态状态。冻结版本将 CSR 事务局部化到 EXU，以减少跨越 IDU、
LSU 和 WBU 的读写控制通路。

\subsection{控制流前端}\label{sec:frontend}

BTB 包含 512 个表项，按 8 个 bank、每 bank 64 项组织为异步 LUTRAM。表项保存类型、tag、压缩目标和
2-bit 饱和方向计数器；更新 shadow 保存较窄的 tag/类型/方向状态。复位阶段逐项清零，在初始化完成前
强制查询 miss。新出现且实际不跳转的条件分支不会驱逐同索引下的有用目标。返回预测由 2 项 RAS
提供，call 压入 \code{pc+4}，标准 return 弹出栈顶。

\begin{signalstable}{BTB 与分支预测器接口信号说明}
query.addr&输入&UInt(32)&组合查询&当前推测 PC\\
query.hit/target&输出&Bool/UInt(32)&组合返回&命中与预测目标\\
query.is*&输出&类型/方向组&随查询返回&区分 branch、JAL、return 及预测方向\\
update.*&输入&地址/目标/类型/方向&注册后写入&用真实执行结果训练表项\\
pred&输出&PredBundle&送入 IFU&给出命中、taken 和下一 PC\\
\end{signalstable}

\begin{figure}[H]\centering\resizebox{0.98\linewidth}{!}{%
\begin{tikzpicture}[
 b/.style={draw,rounded corners=1.2pt,align=center,minimum height=0.82cm,fill=blue!6},
 a/.style={-Stealth,thick},font=\small]
\node[b,minimum width=2.0cm] (pc) at (0,1.9) {推测 PC};
\node[b,minimum width=2.45cm] (split) at (3.0,1.9) {\texttt{pc[16:2]}\\tag / bank / 行号};
\node[b,minimum width=5.6cm,minimum height=1.9cm,fill=green!7] (array) at (8.3,1.9) {
  \textbf{8-bank 并行查询阵列}\\[-0.1em]
  bank 0 \quad bank 1 \quad $\cdots$ \quad bank 6 \quad bank 7\\
  每 bank 64 项，LUTRAM 异步读};
\node[b,minimum width=2.9cm] (pick) at (14.4,1.9) {\texttt{pc[10:8]} 选择 bank\\tag 比较并取出表项};

\node[b,minimum width=3.2cm,fill=orange!8] (decode) at (8.0,-0.55) {类型译码\\JAL / branch / return};
\node[b,minimum width=3.4cm,fill=purple!7] (counter) at (14.0,-0.55) {表项 2-bit 饱和计数器\\最高位给出 branch 方向};
\node[b,minimum width=3.2cm,fill=cyan!7] (ras) at (8.0,-2.8) {2 项返回地址栈\\call 压入，return 弹出};
\node[b,minimum width=4.4cm,fill=gray!10] (next) at (14.0,-2.8) {按类型、方向与 RAS 状态选择\\BTB 目标 / RAS 栈顶 / $pc+4$\\形成 PredBundle};

\draw[a] (pc)--(split);
\draw[a] (split)--(array);
\draw[a] (array)--(pick);
\draw[a] (pick.south west) -- ++(0,-0.42) -| (decode.north);
\draw[a] (pick.south east) -- ++(0,-0.42) -| (counter.north);
\draw[a] (decode)--(ras);
\draw[a] (decode.south east) -- ++(0,-0.36) -| ([xshift=-0.75cm]next.north) -- (next.north);
\draw[a] (counter.south) -- ++(0,-0.42) -| ([xshift=0.75cm]next.north);
\draw[a] (ras)--(next);
\end{tikzpicture}}
\caption{512 项 8-bank BTB 与返回预测结构}\end{figure}


\subsection{存储系统与 D-cache}\label{sec:memory}

D-cache 为 4\,KiB 直接映射结构，由两个 512$\times$32 BRAM bank 组成。tag 使用两个 512$\times$8
分布式存储器；bank 位留在存储地址之外。主数据端是单周期同步读，因此 bank 选择也延迟一拍。
另有每 bank 四个逐字节分布式存储器镜像，提供异步 late-load 数据；该数据必须先跨越 EXU--LSU
寄存边界，再参与受限消费者的运算或比较。

\begin{signalstable}{D-cache 接口信号说明}
queryIndex/tag&输入&10/7 bit&组合 tag 查询&普通 load/store 地址\\
lateQueryIndex&输入&10 bit&异步数据查询&late-load 专用地址\\
hit/readData&输出&Bool/UInt(32)&tag 组合、数据同步&普通缓存命中与主数据\\
lateReadData&输出&UInt(32)&异步读取&逐字节副本数据\\
storeUpdate/*&输入&地址/数据/4-bit mask&store 优先&更新或保守失效缓存行\\
update/*&输入&地址/数据/valid/mask&回填协议&写入较早 load 返回或加速单元更新\\
\end{signalstable}

\begin{figure}[H]\centering\resizebox{\linewidth}{!}{%
\begin{tikzpicture}[
 m/.style={draw,rounded corners=1.2pt,align=center,minimum height=1.0cm,inner sep=5pt},
 a/.style={-Stealth,thick},wa/.style={-Stealth,thick,dashed},font=\small]
\node[m,fill=blue!6,text width=3.35cm] (addr) at (0,2.4)
  {普通访问地址\\tag=[17:11]，bank=[11]\\行=[10:2]};
\node[m,fill=green!7,text width=3.75cm] (tags) at (4.8,2.4)
  {普通 tag LUTRAM\\valid + 7-bit tag，C0 比较};
\node[m,fill=green!7,text width=3.9cm] (main) at (10.1,2.4)
  {双 bank 主数据 BRAM\\每 bank 512$\times$32，C1 同步读};
\node[m,fill=gray!10,text width=3.65cm] (normal) at (15.3,2.4)
  {普通 load 命中\\LSU 对齐、扩展并前递};

\node[m,fill=orange!8,text width=4.05cm] (listaddr) at (0,-0.5)
  {\code{xlistfind} 私有查询\\并行地址 node 与 node+4};
\node[m,fill=orange!8,text width=4.15cm] (replica) at (5.4,-0.5)
  {两组 tag/data 异步副本\\每 bank 四块 512$\times$8 数据 RAM};
\node[m,fill=cyan!8,text width=4.0cm] (walker) at (10.9,-0.5)
  {查找器局部寄存器与状态机\\next/info/data 分阶段解析};
\node[m,fill=cyan!8,text width=3.65cm] (listbus) at (16.0,-0.5)
  {副本 miss\\复用统一数据总线};

\node[m,fill=purple!7,text width=5.35cm] (write) at (3.1,-3.5)
  {注册化统一写口：整字分配 / 窄写失效\\miss 回填 / 硬件加速更新};
\node[m,fill=purple!3,text width=5.1cm] (broadcast) at (10.2,-3.5)
  {同一地址、数据与字节掩码\\更新普通阵列与链表查询副本};

\draw[a] (addr.east) -- (tags.west);
\draw[a] (tags)--node[above]{hit}(main);
\draw[a] (main)--(normal);
\draw[a] (listaddr)--(replica);
\draw[a] (replica)--(walker);
\draw[a] (walker)--node[above]{miss}(listbus);
\draw[a] (write)--(broadcast);
\draw[wa] ([xshift=-1.0cm]broadcast.north) -- ++(0,0.55) -| (replica.south);
\draw[wa] (write.west) -- ++(-3.1,0) -- ++(0,7.7) -| (tags.north);
\draw[wa] (broadcast.east) -- ++(5.6,0) -- ++(0,8.1) -| (main.north);
\end{tikzpicture}}
\caption{总容量 4 KiB 的双 bank D-cache：同步普通加载与链表查找镜像副本}\end{figure}


\subsection{SoC、内存接口与 UART}

CPU 核心通过独立取指与数据 SimpleBus 端口接入 JYD SoC。IROM 基址为 \code{0x80000000}，DRAM
基址为 \code{0x80100000}，MMIO 位于 \code{0x80200000} 区域。FPGA 顶层使用 Vivado BRAM IP
承载指令和数据存储；UART 通过 AXI UARTLite、时钟转换和异步字节 FIFO 连接 CPU 与板级 50\,MHz
外设时钟域，使长时间运行能够留下原始、可机器读取的输出。

\begin{signalstable}{SimpleBus 与 SoC 外设接口信号说明}
req.valid/ready&双向&Bool&请求握手&表示一次请求被接受\\
req.addr&输出&UInt(32)&随请求稳定&统一物理地址\\
wen/data/mask&输出&Bool/32/4 bit&写请求有效&表达 store 数据与字节使能\\
resp.valid/rdata&输入&Bool/UInt(32)&响应握手&返回 load 或取指数据\\
LED/SEG/CNT&输出&设备相关&地址译码&板级显示与计时\\
UART TX/RX&双向&串行字节流&跨时钟 FIFO&输出性能记录并支持交互\\
\end{signalstable}
